%% =====================================================================
%%  T-19-05  The Conscience Attack
%%  -------------------------------------------------------------------
%%  One idea per file. Core is mandatory. Slots in canonical order.
%%  Max 700 words. No layout commands. No filler. New source material
%%  for this entry goes in sources/B6/T-19-05.tex -- never by
%%  rewriting this file (docs/RULES.md R0.1).
%%  =====================================================================
%% @kind: topic
%% @id: T-19-05
%% @title: The Conscience Attack
%% @subtitle: Guilt is priced to the target's conscience --- the more conscientious, the better it works
%% @chapter: c19
%% @order: 50
%% @status: stable
%% @sources: B6
%% @updated: 2026-09-19

\topic{T-19-05}{The Conscience Attack}{Guilt is priced to the target's conscience --- the more conscientious, the better it works}

\begin{core}
Guilt-tripping and shaming are not scattershot; they are calibrated to the target's conscience. B6's clinical finding: manipulators make a tool of their victim's superior scrupulousness, the better to keep them self-doubting, anxious and submissive --- guilt lands hardest on exactly the people most ready to feel it. The attack has two barrels: guilt (you have hurt me) and shame (you are less good than I am).
\end{core}
\begin{mechanism}
The weapon only fires on the armed. A person without conscience shrugs the accusation off; the scrupulous target cannot, because the accusation lands on a live wire --- their own readiness to feel bad, to see the other's side, to blame themselves first. B6's case is Janice and Bill: the drinking husband who \emph{relapses} precisely when she thinks of leaving, so that her exit plans arrive pre-converted into guilt --- and the suicide theatre that seals it, staged after she had been trained to doubt the legitimacy of her own anger. The shaming barrel works beside it: subtle sarcasm and put-downs increase fear and self-doubt, fostering a standing sense of inadequacy that keeps the weaker party deferring. The mechanism's cruelty is structural: the more the target improves themselves --- kinder, more reflective, more willing to own their share --- the better the weapon works. The defence is not to become shameless; it is to price accusations against behaviour and evidence, and to notice that a person who needs you perpetually guilty is not hurt --- they are outfitted.
\end{mechanism}
\begin{conditions}
Strongest in intimate and caretaking relationships, where the target's conscience is engaged daily and the accusation has history to stand on. It weakens in formal settings (records, witnesses), against targets who have done the audit once --- named their own guilt-proneness --- and against genuinely culpable targets, where the guilt is information, not attack. The line between the two is the whole skill.
\end{conditions}
\begin{application}
  \item Audit your guilt reflex: how fast do you move from \emph{they are upset} to \emph{I must have caused it}? That speed is the weapon's fuse.
  \item Price every accusation against the record. Feeling guilty is data about you; the record is data about them.
  \item Separate hurt expressed from hurt staged: real hurt asks for care; staged hurt demands surrender --- watch what the grievance is invoicing for.
  \item Notice the timing correlation. If their collapses cluster around your independence moments, the pattern is the message (T-06-04).
  \item Keep your conscience; audit its use. The goal is a conscience that answers to evidence, not to whoever speaks first.
\end{application}
\begin{feedback}
  \item Working: you can hear an accusation through without settling it instantly --- and most staged ones sound different at the end than at the start.
  \item Working: your anger survives contact; you stop apologising for having limits.
  \item Failing: you are explaining your reasons again, at length, to someone who has stopped listening. The attack achieved its posture.
  \item Failing: you feel guilty about feeling angry. The second barrel hit.
\end{feedback}
\begin{failure}
Over-armouring against guilt is its own failure: a person who treats every accusation as manipulation stops being reachable by genuine hurt, and the people who love them learn to say nothing --- which is its own manipulation-free silence, and its own loss. The doctrine also self-applies: the chronically self-blaming run the attack on themselves without any operator present.
\end{failure}
\begin{countermeasures}
  \item Name the weapon when you see it aimed at someone else: the partner whose every grievance arrives as a bill is running guilt as governance.
  \item Watch the shaming register in workplaces --- the sarcasm that leaves someone perpetually slightly inadequate and therefore deferential. It is a dominance system, not a sense of humour.
  \item For yourself: write down what you actually did before accepting any guilt. Conscience prices feelings; paper prices acts (T-24-06).
  \item If someone collapses whenever you assert a need, say once, plainly, that you notice the pattern. Watch the third collapse's timing, not their tears.
\end{countermeasures}

\sources{B6}
\seesources
\seealso{T-19-01, T-19-02, T-13-02}
